\begin{center} 
\caption{\small Sample attrition} 
\label{tab:attrition} 
\begin{small}
\begin{threeparttable}
\begin{tabular} {@{} l l c c c c c c c c c c c @{}} \toprule \toprule
&&\multicolumn{2}{c}{Grade 1}&&\multicolumn{2}{c}{Grade 2}&&\multicolumn{2}{c}{Grade 3}&&\multicolumn{2}{c}{Grade 4} \\  \cmidrule{3-4} \cmidrule{6-7} \cmidrule{9-10} \cmidrule{12-13}
&&\Longstack{Control\#mean}&\Longstack{Coef.}&&\Longstack{Control\#mean}&\Longstack{Coef.}&&\Longstack{Control\#mean}&\Longstack{Coef.}&&\Longstack{Control\#mean}&\Longstack{Coef.}  \\
&&(1)&(2)&&(4)&(5)&&(7)&(8)&&(10)&(11) \\ \midrule 
\\
&Reading Club Exam &0.75&-0.09&&0.70&-0.03&&0.52&0.17&&0.68&-0.06 \\
&                &         &(0.10)&&         &(0.12)&&         &(0.11)&&         &(0.11) \\
\midrule
\multicolumn{2}{l}{Observations} &547&1,184&&571&1,263&&622&1,243&&509&1,124 \\
\bottomrule
\end{tabular}
\begin{tablenotes}
\item Notes: The table shows the follow-up rate for each group between baseline and endline assessments. The specification used controls for academy cohort dummies. Standard errors are clustered at the academy-grade group level. ***, **, and * indicate significance at 1\%, 5\%, and 10\%.\end{tablenotes}
\end{threeparttable}
\end{small}
\end{center}
\clearpage
